\documentclass[11pt]{article}
\usepackage{amsmath,amssymb,amsthm,mathtools}
\usepackage{microtype}
\usepackage[hidelinks]{hyperref}
\usepackage[nameinlink,capitalize]{cleveref}
\usepackage{booktabs}
\usepackage{enumitem}

\newtheorem{theorem}{Theorem}[section]
\newtheorem{proposition}[theorem]{Proposition}
\newtheorem{lemma}[theorem]{Lemma}
\newtheorem{corollary}[theorem]{Corollary}
\theoremstyle{definition}
\newtheorem{definition}[theorem]{Definition}
\newtheorem{remark}[theorem]{Remark}

\title{Representation-Invariant Geometry of Quantity\\\large From Emergent Carry to Self-Adjoint Realization}
\author{Thiago Massensini}
\date{}

\begin{document}
\maketitle

\begin{abstract}
We study which mathematical structures remain invariant when an exact representation of quantity is changed. Starting from finite local capacity and lossless distinguishability, we derive saturation and carry as transport phenomena, construct a hierarchical depth geometry, identify its natural mass law, and prove a quadratic scaling rigidity whose unique exponent is one half. The resulting real geometry is then developed through reconstruction, positive Gram structure, Hilbert realization, moment and Jacobi geometry, culminating in a real self-adjoint operator constructed downstream of the representation-invariant data.
\end{abstract}

\section{Introduction}
The guiding question is simple: what belongs to quantity itself, and what belongs only to the coordinates used to represent it?

\section{Emergent finite transport}
\section{Hierarchical decomposition of quantity}
\section{Depth mass and quadratic amplitude}
\section{Rigidity of the half-scaling law}
\section{Change of representation and geometric invariance}
\section{Real reconstruction geometry}
\section{Positive Gram and Hilbert realization}
\section{Moment, Hankel, and Jacobi geometry}
\section{Construction of the real operator}
\section{Self-adjoint realization}
\section{Discussion}

\end{document}
